% Chapter 11 worked example. Included by ch11_resources.tex.
\section{Worked Example: Readmission-Risk Prediction: Validation and Causal Evaluation}
\label{sec:comp-ch11-worked-example}

\subsection{Purpose and Status}
This worked example demonstrates how the methods in Chapter~11 can be
translated into a reviewable institutional decision. All numerical inputs are
synthetic unless a source is identified explicitly. They are intended to show the
logic of the analysis, not to report empirical findings or establish a universal
threshold. Local use requires replacement of the assumptions with current,
versioned evidence and approval through the relevant clinical, managerial,
economic, legal, technical, and governance processes.

A hospital has acquired a readmission-risk model intended to direct care-management follow-up. External validation is favourable, but the local population, documentation, social-risk information, and care-management capacity differ from the development setting. 

\subsection{Decision Problem and Comparators}
\textbf{Decision problem.} Is the exact model and intervention sufficiently valid and useful for local use, and does the targeted follow-up pathway cause better outcomes than current discharge planning?

The comparator set is deliberately broader than the existing pathway. This prevents
a new technology from receiving a lower evidentiary threshold than staffing,
workflow, shared-service, conventional digital, or no-change alternatives.

\begin{longtable}{@{}p{0.08\textwidth}p{0.84\textwidth}@{}}
\caption{Comparators for the Chapter 11 worked example}
\label{tab:comp-ch11-comparators}\\
\toprule
\textbf{Option} & \textbf{Comparator or strategy} \\
\midrule
\endfirsthead
\toprule
\textbf{Option} & \textbf{Comparator or strategy} \\
\midrule
\endhead
\bottomrule
\endlastfoot
1 & Current discharge planning \\
2 & Conventional risk checklist \\
3 & Model-guided follow-up with existing capacity \\
4 & Model-guided follow-up with expanded care-management capacity \\
\end{longtable}

\subsection{Model and Decision Structure}
The analytical structure should follow the decision rather than the available
software. The team should map the causal pathway from intervention input to
information, human decision, action, outcome, resource consequence, and review.
The structure should also identify capacity constraints, uncertainty, subgroup
variation, implementation requirements, and events that would change the
recommendation.

A compact quantitative relationship used in this example is:
\begin{equation}
\mathrm{BS}=\frac{1}{n}\sum_{i=1}^{n}(\hat p_i-y_i)^2
\label{eq:comp-ch11-key-relationship}
\end{equation}
The equation is an organizing aid rather than a substitute for disaggregated evidence,
qualitative findings, or mandatory safeguards. Its variables and interpretation must
be defined in the local protocol.

\subsection{Synthetic Parameters and Evidence Sources}
Table~\ref{tab:comp-ch11-parameters} provides the base-case inputs. A
production analysis should add parameter distributions, correlations, structural
alternatives, data dates, system versions, and evidence-confidence ratings.

\begin{longtable}{@{}p{0.32\textwidth}p{0.22\textwidth}p{0.38\textwidth}@{}}
\caption{Synthetic parameters for the Chapter 11 worked example}
\label{tab:comp-ch11-parameters}\\
\toprule
\textbf{Parameter} & \textbf{Base value} & \textbf{Source or interpretation} \\
\midrule
\endfirsthead
\toprule
\textbf{Parameter} & \textbf{Base value} & \textbf{Source or interpretation} \\
\midrule
\endhead
\bottomrule
\endlastfoot
Annual eligible discharges & 18,000 & Hospital activity \\
30-day readmission prevalence & 14\% & Baseline cohort \\
Available care-management slots & 65/week & Service capacity \\
Model sensitivity at proposed threshold & 0.76 & Local silent evaluation \\
Positive predictive value & 0.31 & Local silent evaluation \\
Calibration slope & 0.82 & Local validation \\
Expected follow-up effect & Relative risk 0.88 & Evidence synthesis; causal uncertainty remains \\
\end{longtable}

\subsection{Step-by-Step Analysis}
The sequence below is intended to make the analysis reproducible and to ensure
that evidence, economics, governance, and implementation develop together.

\begin{longtable}{@{}p{0.07\textwidth}p{0.55\textwidth}p{0.30\textwidth}@{}}
\caption{Analytical sequence}
\label{tab:comp-ch11-analysis-steps}\\
\toprule
\textbf{Step} & \textbf{Required work} & \textbf{Decision record} \\
\midrule
\endfirsthead
\toprule
\textbf{Step} & \textbf{Required work} & \textbf{Decision record} \\
\midrule
\endhead
\bottomrule
\endlastfoot
1 & Define the intended use, threshold, action, capacity, and exact system version. & Evidence and responsible owner to be recorded \\
2 & Assess data design, leakage, transportability, discrimination, calibration, threshold-specific utility, and subgroup performance. & Evidence and responsible owner to be recorded \\
3 & Conduct silent evaluation before clinical dependence. & Evidence and responsible owner to be recorded \\
4 & Separate model validation from causal evaluation of the follow-up intervention. & Evidence and responsible owner to be recorded \\
5 & Use an appropriate randomized, quasi-experimental, or target-trial design to estimate service effects. & Evidence and responsible owner to be recorded \\
6 & Combine technical, clinical, implementation, qualitative, economic, governance, and post-deployment evidence in a decision-ready record. & Evidence and responsible owner to be recorded \\
\end{longtable}

\subsection{Illustrative Outputs}
The outputs in Table~\ref{tab:comp-ch11-outputs} show how technical,
clinical, operational, economic, equity, and governance findings should be kept
visible rather than compressed into a single favourable or unfavourable score.

\begin{longtable}{@{}p{0.30\textwidth}p{0.24\textwidth}p{0.38\textwidth}@{}}
\caption{Illustrative model and decision outputs}
\label{tab:comp-ch11-outputs}\\
\toprule
\textbf{Output} & \textbf{Result} & \textbf{Interpretation} \\
\midrule
\endfirsthead
\toprule
\textbf{Output} & \textbf{Result} & \textbf{Interpretation} \\
\midrule
\endhead
\bottomrule
\endlastfoot
Local discrimination & Acceptable & Not sufficient for adoption \\
Calibration & Overprediction in higher-risk groups & Requires recalibration or threshold review \\
Weekly patients flagged & Approximately 120 & Exceeds current care-management capacity \\
Causal effect of follow-up & Uncertain locally & Requires controlled evaluation \\
Preferred decision & Recalibrate and conduct capacity-constrained pilot & No unrestricted implementation \\
\end{longtable}

\subsection{Sensitivity, Uncertainty, and Subgroups}
At minimum, the completed analysis should vary the parameters most likely to
change the decision, test at least one credible alternative model structure, and
report subgroup results separately where performance, access, response, burden,
or opportunity cost may differ. Deterministic analysis should identify thresholds;
probabilistic analysis should be used where credible parameter distributions and
correlations are available. Scenario analysis is preferable when structural or
future uncertainty cannot be represented defensibly by one probability model.

The record should distinguish uncertainty that can be reduced through evidence,
uncertainty that can be managed through safeguards, and uncertainty that remains
too consequential to accept. Missing evidence should not be represented as an
average or neutral value without justification.

\subsection{Interpretation and Recommendation}
Complete recalibration and a prospective, capacity-aware evaluation of the model-plus-follow-up intervention. The model should not be judged effective merely because it predicts readmission; the evaluation must show that targeted action improves outcomes relative to realistic alternatives.

The recommendation is conditional rather than permanent. It applies only to the
specified population, setting, intervention configuration, evidence cut-off, and
version. The following conditions should be incorporated into the approval,
contract, implementation, monitoring, or renewal record:

\begin{itemize}
 \item silent-mode evidence and subgroup calibration are acceptable.
 \item the threshold fits available care-management capacity.
 \item the causal evaluation separates model performance from intervention effect.
 \item implementation costs and workload are measured prospectively.
 \item the evidence record states attribution limits and conditions for continuation.
\end{itemize}

\subsection{Decision Statement}
\begin{quote}
For the specified decision problem and comparator set, the institution should
\emph{[proceed / proceed with conditions / redesign / pilot / restrict / defer /
replace / retire]}. The decision applies to \emph{[population, setting, version,
and evidence period]}, is owned by \emph{[accountable authority]}, and will be
reviewed on \emph{[date]} or earlier if \emph{[trigger events]} occur. The
evidence, assumptions, unresolved uncertainty, mandatory safeguards, monitoring
thresholds, and exit arrangements are recorded in the accompanying dossier.
\end{quote}
